\documentclass[a4paper,11pt]{ltjsarticle}
\usepackage[haranoaji]{luatexja-preset}
\usepackage[margin=24mm]{geometry}
\usepackage{fix-cm}
\usepackage{amsmath,amssymb,amsthm,mathtools,booktabs,array}
\usepackage[unicode,hidelinks]{hyperref}
\usepackage{xurl}
\setlength{\emergencystretch}{3em}
\newtheorem{theorem}{定理}[section]
\newtheorem{proposition}[theorem]{命題}
\newtheorem{lemma}[theorem]{補題}
\newtheorem{corollary}[theorem]{系}
\theoremstyle{remark}\newtheorem{remark}[theorem]{注意}
\newcommand{\PP}{\mathbf P}\newcommand{\CC}{\mathbf C}\newcommand{\QQ}{\mathbf Q}
\newcommand{\Jac}{\operatorname{Jac}}\newcommand{\gr}{\operatorname{gr}}
\newcommand{\Spec}{\operatorname{Spec}}\newcommand{\rank}{\operatorname{rank}}
\newcommand{\Span}{\operatorname{Span}}
\title{完全等方Log点におけるPoisson接錐の非被約性\\
\large 二次障害の完全計算・重複度一・混合障害の消去}
\author{}
\date{2026年11月10日\quad v13・研究原稿}
\begin{document}\maketitle
\begin{abstract}
複素射影四空間のPoissonスキームの一つの階数二点について、最小局所埋込みの二次初期方程式を完全に計算する。
接空間は30次元で、独立な二次方程式は7本である。
その二次障害スキームの被約台は次元28、26、24の三既約成分を持ち、
非零の二次式が平方零元を与える。この二次式は真の接錐でも消えないため、
\emph{真の接錐自体が非被約}である。一方、最大次元の接錐成分は単一の28次元線形空間で、
その一般点での長さは一である。従って局所次元は28、Hilbert--Samuel重複度は一となる。これを用いて28次元の階数二対数型成分が交点でも滑らかであることを導き、25次元の階数四長鎖成分との交叉次元を20以上24以下と評価する。
さらに、固定二次超曲面上の延長ファイバーに現れた混合関係式$ac=0$は、
全周囲Poisson変形では明示的な二次補正により消去できることを示す。
二次初期イデアルと全初期イデアル、接錐の非被約性と元の局所環の非被約性は区別する。
\end{abstract}

\section{問題設定と先行研究との関係}
非零Poissonテンソルを尺度で射影化したスキームを
$\mathcal P\subset\PP H^0(\PP^4,\wedge^2T_{\PP^4})$とする。
$\mathrm{PGL}_5$による商は取らない。
本稿の対象は、階数二の対数型構造と階数四の構造が交わる具体的な点である。
この点の固定二次超曲面に関する延長スキーム、局所次元、既知の階数四成分への帰属は
同伴原稿\cite{Incidence}で研究した。本稿は接錐のscheme構造に焦点を当て、
局所次元の計算に必要な係数証明書も独立に再構成する。

Matviichuk--Pym--Schedler\cite{MPS}は、正規交叉極因子を持つ
log symplectic構造の近傍の変形空間を研究する。
ここで扱う中心点は一般階数二であり、一般階数四のlog symplectic点そのものではない。
従って同論文の局所模型をそのまま本点に適用するのではなく、
全周囲Jacobi方程式から障害写像を直接計算する。
先行研究全体に対する新規性・優先権の確認は未完了であり、
以下の具体的恒等式の証明と、新規性の主張を混同しない。
重複度の可換環論には\cite{StacksMult}の加法性・結合公式を用いる。

\section{模型と接空間の座標}
同次座標を$X=(u_0,u_1,z,w_0,w_1)$、体積形式を$\Omega$、Euler場を$E$とする。
\begin{equation}\label{eq:model}
 q=u_0^2+2u_1w_0+2zw_1,\qquad
 \xi=u_0\partial_{u_0}+u_1\partial_{u_1}-2z\partial_z,
 \qquad
 \iota_{P_0}\Omega=dq\wedge\iota_\xi(du_0\wedge du_1\wedge dz).
\end{equation}
この$P_0$は非零の階数二Poisson双ベクトルである。
$p=[P_0]$、$R=\mathcal O_{\mathcal P,p}$、極大イデアルを$\mathfrak m$とする。
斉次二次双ベクトル$B$のJacobiatorを
\[
 J(B)_{ijk}=\sum_{\ell=0}^4
 (B_{i\ell}\partial_\ell B_{jk}+B_{j\ell}\partial_\ell B_{ki}
                              +B_{k\ell}\partial_\ell B_{ij})
\]
と書く。射影Poisson方程式は$\mathcal F(B)=E\wedge J(B)=0$である。
二次双ベクトルには150係数があり、$\mathcal F$には350個の四次係数がある。
Eulerゲージ$E\wedge$一次場の像は24次元である。

以下に接空間の基底を固定する。添字は$X=(X_0,\ldots,X_4)$の順で、全て0始まりとする。
$D_{ij}=[X_j\partial_{X_i},P_0]$について、$5i+j$が
\[
 (0,1,2,3,4,5,6,7,8,9,10,11,13,14,16,17,19,20,22,23)
\]
である20元をこの順に取る。続いて式\eqref{eq:model}で$q$を
\[
 u_0w_1,\quad w_0^2,\quad w_0w_1,\quad w_1^2
\]
の方向に変える4元、$\xi$を
\[
 u_1\partial_{u_0},\quad u_0\partial_{u_1},\quad
 u_0\partial_{u_0}-z\partial_z
\]
の方向に変える3元を取る。最後の3元を次で定める。
$G=\operatorname{Hess}(q)/2$とし、$A$を$A^tG+GA=0$を満たす行列とするとき
\begin{equation}\label{eq:K}
 K_A=\frac q2AG^{-1}+\frac12E\wedge(AX).
\end{equation}
\[
 A_1X=u_1\partial_{u_0}-u_0\partial_{w_0},\quad
 A_2X=z\partial_{w_0}-u_1\partial_{w_1},\quad
 A_3X=u_1\partial_{u_1}-z\partial_z-w_0\partial_{w_0}+w_1\partial_{w_1}
\]
に対応する$K_1,K_2,K_3$を加える。場とその成分ベクトルはここで同一視する。
この30元を$B_0,\ldots,B_{29}$、対応する接座標を$t_0,\ldots,t_{29}$とする。
\begin{equation}\label{eq:coords}
 s_i=t_{20+i}\ (0\le i\le6),\qquad a=t_{27},\ b=t_{28},\ c=t_{29}.
\end{equation}
最初の20個の軌道方向と$b$は、下記の二次関係に現れない。

\section{全二次初期イデアル}
$S=\CC[t_0,\ldots,t_{29}]$とする。
$\gr_{\mathfrak m}R=S/\mathcal I$という接空間からの表示において、
$\mathcal I$の次数二部分が生成するイデアルを$I_2$と書く。
$I_2$と全初期イデアル$\mathcal I$は一般には異なる。

\begin{theorem}[二次方程式の完全性]\label{thm:quadrics}
$\dim T_p\mathcal P=30$であり、上記の座標において
\begin{equation}\label{eq:seven}
\boxed{\quad
 I_2=(s_0c-s_2a,\ s_1c,\ c(s_2+c),\ s_3c,
             \ s_4c+s_6a,\ s_5a,\ s_5c).\quad}
\end{equation}
特に$\dim_\CC\mathcal I_2=7$である。これは全接錐が二次式だけで定義されるという主張ではない。
\end{theorem}
\begin{proof}
$M=d\mathcal F_{P_0}$を全350係数、全150二次双ベクトル方向について作る。
有理数体上の階数は95である。
上述の30元は全て$M$の核に入り、Eulerゲージ24元および$P_0$を加えた行列は階数55を持つ。
従って尺度も割った接空間で30元は基底をなす。

$M$の像に写る95個の補空間座標を選び、線形項が独立な95個の方程式を形式的に解く。
残る最小局所表示は30変数で、全定義方程式は二次以上から始まる。
その二次部分は、$M$の全左核$\Lambda M=0$に対する
\begin{equation}\label{eq:obmap}
 \Lambda\!\left(\mathcal F\!\left(\sum_{i=0}^{29}t_iB_i\right)\right)
\end{equation}
が張る空間に正確に等しい。高次の係数を掛ける操作や方程式間の消去は、
この次数二部分に新しい独立な関係を加えない。

この有限係数計算を明示的に特定する。
添字対は辞書式、次数$d$単項式は$0\le i_1\le\cdots\le i_d\le4$の辞書式で並べる。
$M^t$の被約行階段形のピボット列を$p_0,\ldots,p_{94}$とする。
各非ピボット列$j$に対し
$\lambda_j=e_j-\sum_k\operatorname{rref}(M^t)_{kj}e_{p_k}$と置けば、
これら255元は全左核の基底である。
30個の対角項と435個の異なる基底対に式\eqref{eq:obmap}を評価する。
得られる$255\times465$行列の階数は7で、非零の被約行階段形は式\eqref{eq:seven}である。
非ピボット列を昇順に番号付けたとき、次の7行と倍率で正規化された7式を直接回収できる。
\[
\begin{array}{c|rrrrrrr}
 \text{式の順番}&1&2&3&4&5&6&7\\\hline
 \text{左核の行番号}&140&52&47&145&110&95&88\\
 \text{倍率}&-1/12&1/9&-10/9&-2/33&1/2&-1/4&1/2
\end{array}
\]
全350係数、全150方向、全465組の検算と左核汎関数の疎な全表示を、付属する
\path{code/subvariety_incidence/check_isotropic_quadratic_cone.py}
および対応するJSON証明書で、外部入力なしに再生成する。
従ってこれは標本点や一部の接方向だけについての計算ではない。
\end{proof}

\section{二次障害スキームと真の接錐の非被約性}
\begin{theorem}[被約台と平方零部分]\label{thm:radical}
$h=a(s_2+c)$と置く。このとき
\begin{align}
 \sqrt{I_2}=I_2+(h)&=P_{28}\cap P_{26}\cap P_{24},\label{eq:radical}\\
 P_{28}&=(a,c),\nonumber\\
 P_{26}&=(c,s_2,s_6,s_5),\nonumber\\
 P_{24}&=(s_1,s_3,s_5,s_2+c,s_0+a,s_4c+s_6a).\nonumber
\end{align}
添字は対応する$S$上の既約成分の次元である。
$S/I_2$のnilradicalは$\bar h$により生成され、
\begin{equation}\label{eq:nil}
 \bar h\ne0,\quad\bar h^2=0,\qquad
 I_2:h=(s_1,s_2,s_3,s_5,s_6,c).
\end{equation}
さらに、\emph{真の接錐}$\Spec\gr_{\mathfrak m}R$にも$h$の像は非零の平方零元として残る。
従って真の接錐は非被約である。
\end{theorem}
\begin{proof}
三つの表示イデアルは素である。$P_{24}$の場合は5個の独立な一次式を除いた後に
既約二次式$s_4c+s_6a$が残るからである。
三素イデアルの交叉の生成元は
\[
\begin{split}
 &as_6(s_0+a),\ c(s_0+a),\ as_6s_1,\ cs_1,\ a(s_2+c),\ c(s_2+c),\\
 &as_6s_3,\ cs_3,\ s_4c+s_6a,\ as_5,\ cs_5.
\end{split}
\]
これは$I_2+(h)$と一致する。
例えば$P_{28}\cap P_{26}=(c,as_2,as_6,as_5)$から残りの交叉を計算できる。
各包含は上の生成元への直接代入でも確認でき、付属コードは両方向のGr\"obner剰余を検査する。

第1式を$g_1=s_0c-s_2a$、第3式を$g_3=c(s_2+c)$と書けば、恒等的に
\begin{equation}\label{eq:h-square}
 h=c(s_0+a)-g_1,\qquad
 h^2=a(s_0+a)g_3-hg_1.
\end{equation}
よって$h^2\in I_2$である。
一方、$h$の単項式$ac$の係数は1であり、式\eqref{eq:seven}のどの生成元にも$ac$は現れない。
次数二により$h\notin I_2$となる。
これと式\eqref{eq:radical}からnilradicalの記述が従う。

式\eqref{eq:h-square}と7式を使うと$s_1,s_2,s_3,s_5,s_6,c$は全て$\bar h$を殺す。
逆包含は$\CC(s_0,s_4,a)$上の双対数への写像
\[
 c\mapsto\varepsilon,\quad s_2\mapsto(s_0/a)\varepsilon,\quad
 s_6\mapsto-(s_4/a)\varepsilon,\quad s_1,s_3,s_5\mapsto0,
 \qquad\varepsilon^2=0
\]
で確認できる。ここで$h\mapsto(s_0+a)\varepsilon$であり、
$s_0,s_4,a$および残る自由変数の非零多項式はこれを殺さない。

最後に定理\ref{thm:quadrics}は、単に7式が真の接錐を消すというだけでなく、
\emph{真の初期イデアルの次数二部分全体}を決定している。
従って$h\notin\mathcal I_2$であり、その像は$S/\mathcal I$でも非零である。
式\eqref{eq:h-square}からその平方は零である。
高次の初期方程式が追加されても、この次数二の非零性は失われない。
\end{proof}

\begin{remark}[区別すべき三つのスキーム]
$\Spec(S/I_2)$は二次障害スキームであり、真の接錐はその閉部分スキームである。
次元26、24の二次障害成分が、そのまま真の接錐や元のPoissonスキームの成分になるとは限らない。
接錐の非被約性は証明したが、元の局所環$R$の非被約性は未証明である。
真の接錐に残る$\bar h$の台は、式\eqref{eq:nil}の6一次式で定まる次元24の空間に含まれる。
その台の正確な次元は本稿からは決まらない。
\end{remark}

\section{局所次元と重複度一}
\begin{lemma}[独立な28次元下界]\label{lem:lower}
$p$を含む階数二対数型Poisson族の閉包は28次元である。
\end{lemma}
\begin{proof}
$\omega_0=\iota_E\iota_{P_0}\Omega$は
\[
 \omega_0=qu_0u_1z\,d\log(u_0/u_1)\wedge d\log(q^3/(u_1^4z^2))
\]
と書ける。留数は$(1,-1,0,0)$と$(0,-4,-2,3)$で、どちらも$(1,1,1,2)$と直交する。
一般の対数型源は、三一次式$f_0,f_1,f_2$と一二次式$f_3$、
$c\in\PP\wedge^2\ker(1,1,1,2)\simeq\PP^2$から
\[
 \omega(f,c)=\sum_{i<j}c_{ij}\prod_{k\ne i,j}f_k\,df_i\wedge df_j
\]
を作るもので、源の次元は$12+14+2=28$である。
三次元空間の二重外積の元は分解可能である。二つの閉対数一形式の共通核は可積分であり、対応する分解可能双ベクトルはPoissonである。

像の次元に関する下界を正確な微分の小行列式で確認する。
$f_0=u_0,f_1=u_1,f_2=z$、
\[
\begin{split}
f_3={}&u_0^2+u_1^2+z^2+w_0^2+w_1^2+u_0u_1+u_0w_0+u_0w_1\\
     &+2u_1w_0+u_1w_1+zw_0+2zw_1+w_0w_1
\end{split}
\]
とする。$r_1=(1,-1,0,0)$、$r_2=(1,0,-1,0)$、$r_3=(2,0,0,-1)$により
$c=r_1\wedge r_2+2r_1\wedge r_3+5r_2\wedge r_3$と置く。
$f_0,f_1,f_2$のそれぞれの主係数を固定した12変数、$f_3$の$w_1^2$係数を固定した14変数、
$c$の第1係数を固定した2変数で微分する。
350個の三次二形式係数を前節と同じ辞書式順で並べた微分行列の行
\[
\begin{split}
(2,3,4,6,7,8,9,10,11,16,19,20,21,25,26,27,28,29,\\
30,38,39,42,43,45,46,52,53,55)
\end{split}
\]
の小行列式は$-57177414000000\ne0$である。
基点$\omega(f,c)$の係数列も加えると階数29なので、射影化しても微分階数28を持つ。
付属コードはこの小行列式と$p$の対数型表示を独立に再計算する。
従って像の閉包は28次元であり、表示式から$p$を含む。
\end{proof}

\begin{theorem}[最大次元接錐とHilbert--Samuel重複度]\label{thm:mult}
\[
 \dim R=28,\qquad (C_p\mathcal P)_{\mathrm{top,red}}=V(a,c)\simeq\mathbf A^{28},
 \qquad e(R)=1.
\]
$\mathrm{top,red}$は最大次元既約成分の被約な台を意味する。
\end{theorem}
\begin{proof}
真の接錐は$\Spec(S/I_2)$に含まれ、後者の最大次元は28である。
補題\ref{lem:lower}により逆の下界も成立し、$\dim R=28$となる。
定理\ref{thm:radical}の三素イデアルのうち28次元なのは$P_{28}=(a,c)$のみだから、
真の接錐の28次元の被約台はこの線形空間に一致する。
この素イデアルで局所化すると$s_1,s_5$は単元であり、$s_1c$と$s_5a$から
$(I_2)_{P_{28}}=(P_{28})_{P_{28}}$を得る。
従って真の接錐のこの一般点での長さは一である。
局所環の重複度は接錐の次数に等しく、最大次元成分への結合公式
\cite[Lemma 43.15.3]{StacksMult}により、長さ一と線形空間の次数一の積、すなわち一となる。
\end{proof}

\begin{remark}[元の非Cohen--Macaulay性との関係]
同伴原稿\cite{Incidence}で示した25次元の階数四長鎖成分への帰属と、
28次元対数型成分の成分性を用いると、$R$は非等次元で、深さは25以下である。
その成分性入力は本稿の二次計算だけの帰結ではない。
本稿の$e(R)=1$はそれと矛盾せず、低次元成分や低次元に支持されるnilpotent構造は
最高次のHilbert--Samuel係数に寄与しない。
\end{remark}


\begin{theorem}[最大次元の対数型成分は交点でも滑らか]
\label{thm:log-smooth-at-intersection}
$Z_{\mathrm{log}}=\overline{\operatorname{Log}(1,1,1,2)}$を
補題\ref{lem:lower}の28次元階数二対数型既約閉包とする。
このとき$p\in Z_{\mathrm{log}}$は$Z_{\mathrm{log}}$の滑らかな点であり、
\begin{equation}\label{eq:smooth-log-branch}
 T_pZ_{\mathrm{log}}=V(a,c)\subset T_p\mathcal P\simeq\mathbf A^{30},
 \qquad
 \widehat{\mathcal O}_{Z_{\mathrm{log}},p}
 \cong \CC[[t_0,\ldots,t_{26},b]].
\end{equation}
右辺の同型は、表示した28個の接空間座標の制限により誘導される。
従って、元のPoissonスキームの28次元の成分は$p$で
重複度一の正則な形式的枝である。
\end{theorem}
\begin{proof}
$R=\mathcal O_{\mathcal P,p}$とし、
$Z_{\mathrm{log}}$の極小素イデアルを$\mathfrak p_{\mathrm{log}}$とする。
定理\ref{thm:mult}より$\dim R=28$かつ$e(R)=1$である。
Hilbert--Samuel重複度の結合公式\cite{StacksMult}を
最大次元の極小素イデアルに適用すると
\[
 1=e(R)=\sum_{\substack{\mathfrak p\in\operatorname{Min}(R)\\
                  \dim(R/\mathfrak p)=28}}
 \ell_{R_{\mathfrak p}}(R_{\mathfrak p})\,e(R/\mathfrak p).
\]
各項は正の整数であり、$Z_{\mathrm{log}}$は28次元で$p$を通るから、
唯一の項は$\mathfrak p_{\mathrm{log}}$で
\[
 e(\mathcal O_{Z_{\mathrm{log}},p})=1,\qquad
 \ell_{R_{\mathfrak p_{\mathrm{log}}}}(R_{\mathfrak p_{\mathrm{log}}})=1.
\]
$\mathcal O_{Z_{\mathrm{log}},p}$は複素数体上有限型の
既約多様体の局所整域であり、excellentかつformally unmixedである。
Nagataの重複度一判定\cite[Theorem 40.6]{NagataLocal}
からこれは28次元の正則局所環となる。

よって$Z_{\mathrm{log}}$の接錐は28次元の線形空間
$T_pZ_{\mathrm{log}}$であり、$C_p\mathcal P$に閉埋込みされる。
定理\ref{thm:radical}、\ref{thm:mult}より、
$C_p\mathcal P$の被約台の28次元既約成分は
$V(a,c)$のみであり、他はこれより低次元である。
従って同じ次元の線形空間
$T_pZ_{\mathrm{log}}$は$V(a,c)$に一致する。
この平面上の$t_0,\ldots,t_{26},b=t_{28}$は
正則局所環の余接空間の基底なので、
完備化の形式座標を与える。
\end{proof}

\begin{corollary}[二成分の交叉の次元範囲]
\label{cor:log-long-intersection-orbit}
同伴原稿で同定した25次元の一般階数四既約成分
$Z_{\mathrm{long}}$に対し、
\[
 \boxed{\quad
 20\le\dim_p(Z_{\mathrm{log}}\cap Z_{\mathrm{long}})\le24.\quad}
\]
両成分は$p$を通る20次元の$\mathrm{PGL}_5$軌道を含む。
\end{corollary}
\begin{proof}
同次斉次二次双ベクトル$P_0$に
$\mathfrak{gl}_5$の25個の行列単位を作用させると、
全150係数への微分の階数は$\QQ$上で21となる。
$P_0$はこの像に属し、$P_0$のスカラー倍を割った
射影軌道の次元は$21-1=20$である。
その全係数計算は付属の
\texttt{check\_isotropic\_orbit20.py}で独立に再現できる。
$Z_{\mathrm{log}}$と$Z_{\mathrm{long}}$は
$\mathrm{PGL}_5$不変な閉既約集合で、両方$p$を含む。
従って両者の交叉は20次元の軌道を含み、下界を得る。
一方、$Z_{\mathrm{long}}$の一般点は階数四だが
$Z_{\mathrm{log}}$は全点で階数二以下だから、
$Z_{\mathrm{long}}\not\subset Z_{\mathrm{log}}$である。
よって交叉は25次元の既約集合の真の閉部分集合で、
$p$での局所次元は高々24となる。
\end{proof}

\begin{remark}
$\mathcal O_{Z_{\mathrm{log}},p}$の滑らかさは、
$R=\mathcal O_{\mathcal P,p}$の滑らかさや被約性を意味しない。
$R$には別の25次元成分が通り、
真の接錐には非零平方零元が残る。
また、交叉の正確な次元やそのscheme構造も未決定である。
\end{remark}

\section{混合\texorpdfstring{$ac$}{ac}障害の周囲での消去}
固定された$Q=V(q)$と$P_0|_Q$への延長ファイバーは、同伴原稿で
$\Spec\CC[a,b,c]/(ac,c^2)$と計算されている。
$c^2$の障害は全周囲でも残るが、$ac$には異なる挙動がある。

\begin{proposition}[明示的な混合補正]\label{prop:mixed}
二次双ベクトル$R_{ac}$の非零係数を座標順$(u_0,u_1,z,w_0,w_1)$で
\[
\begin{aligned}
 (R_{ac})_{02}&=-\tfrac12(u_0^2-u_1w_0+2zw_1),&
 (R_{ac})_{03}&=-\tfrac32w_0w_1,\\
 (R_{ac})_{12}&=-\tfrac34u_0u_1,&
 (R_{ac})_{13}&=\tfrac34u_0w_1,&
 (R_{ac})_{23}&=\tfrac34u_0w_0
\end{aligned}
\]
とし、反対称性で残りを定める。このとき
\[
 P_0+aK_1+cK_3+acR_{ac}
\]
は係数環$A=\CC[a,c]/\bigl((a,c)^3+(c^2)\bigr)$上で射影Poissonテンソルを定める。
\end{proposition}
\begin{proof}
$M=d\mathcal F_{P_0}$、
$\mathcal F(B,C)=\mathcal F(B+C)-\mathcal F(B)-\mathcal F(C)$と書く。
全350係数についての正確な恒等式
\[
 M(R_{ac})=-\mathcal F(K_1,K_3)
\]
が成立する。さらに$M(K_1)=M(K_3)=0$および$\mathcal F(K_1)=0$である。
従って全二次までの残差は$c^2\mathcal F(K_3)$だけとなり、$A$上では零である。
これらの恒等式は同じ独立producerで全係数を照合した。
これは$ac$項の障害の消去であり、$c$方向そのものの二次持上げや全次数の変形の存在ではない。
\end{proof}

\section{再現・依存関係・未解決箇所}
計算環はPython 3、SymPy 1.14.0、有理数体である。
前のproducerをimportせず、全行列・全二次障害・対数型源の小行列式・混合補正を再構成する。
\begin{verbatim}
python code/subvariety_incidence/check_isotropic_quadratic_cone.py \
  --output work/isotropic_quadratic_cone.json \
  --check research/subvariety_incidence/certificates/\
isotropic_quadratic_cone.json
latexmk papers/isotropic_log_tangent_cone_ja.tex
\end{verbatim}
検算済みの結果と、本稿からは決まらない部分を分ける。
全二次初期イデアル、真の接錐の非被約性、局所次元28、重複度一、混合$ac$補正は証明済みである。
全初期イデアル、接錐の低次元成分の正確な分解、$R$そのものの被約性、
$p$を通る全既約成分の列挙および全次数のPoisson変形は未解決である。
二次障害スキームの次元26、24を新しいPoisson成分として数えてはならない。

\begin{thebibliography}{9}
\bibitem{NagataLocal} M. Nagata, \emph{Local Rings},
Interscience Tracts in Pure and Applied Mathematics, No.~13,
Interscience Publishers, 1962, Theorem 40.6.
\bibitem{MPS} M. Matviichuk, B. Pym and T. Schedler,
\emph{A local Torelli theorem for log symplectic manifolds},
arXiv:2010.08692v2 (2021), \url{https://arxiv.org/abs/2010.08692v2}.
正規交叉極因子を持つlog symplectic領域の局所変形理論を背景として参照。
本点への直接適用は行わない。
\bibitem{StacksMult} The Stacks Project Authors,
\emph{The Stacks Project}, Section 43.15, especially Lemmas 43.15.2--43.15.3,
\url{https://stacks.math.columbia.edu/tag/0AZU}.
\bibitem{Incidence} 本研究の同伴原稿、
\emph{Poisson部分多様体から見る射影四空間の成分分類}、2026年11月10日、現行研究原稿。
\path{papers/poisson_subvariety_incidence_ja.tex}、共有証明
\path{papers/quadric_log_component_intersection.tex}、
\path{papers/rank_four_long_component_membership.tex}。
外部公表・査読済みの文献としては扱わない。
\end{thebibliography}
\end{document}
